
\subsubsection{5.1 Discriminative Degradation Exists (H-E1)}

Table 1 presents AUROC comparisons between base and instruction-tuned models.

\textbf{Table 1: AUROC for Margin-Based Correctness Prediction}

\begin{table*}[htbp]
\centering
\resizebox{\dimexpr 2\columnwidth+\columnsep\relax}{!}{%
\begin{tabular}{lllll}
\toprule
Family & Base AUROC & Instruct AUROC & Delta & 95\% CI \\
\midrule
Qwen & 0.8298 & 0.8076 & +0.0222 & [0.0074, 0.0370] \\
Mistral & 0.7797 & 0.7413 & +0.0385 & [0.0238, 0.0531] \\
Mean & - & - & +0.0303 & - \\
\bottomrule
\end{tabular}
}
\end{table*}

Both model families show statistically significant AUROC decreases in instruction-tuned models, with confidence intervals excluding zero. The mean degradation is 3.03 percentage points.

Additional context: Base model accuracies were 0.7139 (Qwen) and 0.5875 (Mistral); instruction-tuned accuracies were 0.7045 (Qwen) and 0.5772 (Mistral). The AUROC degradation is not explained by accuracy differences, as AUROC measures rank-ordering independent of accuracy level.

\subsubsection{5.2 Margin Inflation Mechanism (H-M1)}

Table 2 presents conditional margin statistics.

\textbf{Table 2: Expected Margin Conditional on Correctness}

\begin{table*}[htbp]
\centering
\resizebox{\dimexpr 2\columnwidth+\columnsep\relax}{!}{%
\begin{tabular}{llllll}
\toprule
Family & $E[m \mid \text{incorrect}]_{\text{base}}$ & $E[m \mid \text{incorrect}]_{\text{inst}}$ & Inflation Ratio & Cohen's d & p-value \\
\midrule
Qwen & 0.9597 & 2.9327 & 3.06x & 1.01 & $<$0.0001 \\
Mistral & 0.4682 & 7.8606 & 16.79x & 1.85 & $<$0.0001 \\
\bottomrule
\end{tabular}
}
\end{table*}

For context, expected margins conditional on correct predictions also increased: Qwen from 3.6050 (base) to 7.1491 (instruct), and Mistral from 1.6236 (base) to 13.3810 (instruct). However, the inflation is proportionally larger for incorrect predictions, particularly in Mistral where incorrect-margin inflation (16.79x) substantially exceeds correct-margin inflation (~8x).

The large effect sizes (Cohen's d > 1.0) indicate this is not a subtle statistical effect but a substantial behavioral difference.

\subsubsection{5.3 Percentile-Normalized Slope Attenuation (H-M2)}

Table 3 presents logistic regression slopes after percentile normalization.

\textbf{Table 3: Percentile-Normalized Beta Coefficients}

\begin{table*}[htbp]
\centering
\resizebox{\dimexpr 2\columnwidth+\columnsep\relax}{!}{%
\begin{tabular}{llllll}
\toprule
Family & $\beta _base$ & $\beta _instruct$ & Delta $\beta$ & 95\% CI & p-value \\
\midrule
Qwen & 2.2222 & 1.4661 & 0.7581 & [0.6626, 0.8563] & <0.0001 \\
Mistral & 1.5579 & 0.9305 & 0.6284 & [0.5570, 0.7005] & <0.0001 \\
\bottomrule
\end{tabular}
}
\end{table*}

Slope attenuation persists after removing scale differences through percentile normalization. The slopes decrease by 34\% (Qwen) and 40\% (Mistral), indicating that the confidence-correctness relationship is weakened in shape, not merely rescaled.

\subsubsection{5.4 Brier Decomposition Confirms Geometric Distortion (H-M3)}

Table 4 presents Brier score decomposition results.

\textbf{Table 4: Brier Score Decomposition}

\begin{table*}[htbp]
\centering
\resizebox{\dimexpr 2\columnwidth+\columnsep\relax}{!}{%
\begin{tabular}{lllll}
\toprule
Family & Component & Base & Instruct & Delta \\
\midrule
Qwen & Brier Score & 0.1647 & 0.2263 & +0.0616 \\
Qwen & Reliability & 0.0173 & 0.0547 & +0.0374 \\
Qwen & Refinement & 0.0559 & 0.0343 & -0.0216 \\
Qwen & Uncertainty & 0.2042 & 0.2082 & +0.0040 \\
Mistral & Brier Score & 0.2023 & 0.3838 & +0.1814 \\
Mistral & Reliability & 0.0185 & 0.1507 & +0.1321 \\
Mistral & Refinement & 0.0580 & 0.0093 & -0.0487 \\
Mistral & Uncertainty & 0.2424 & 0.2440 & +0.0017 \\
\bottomrule
\end{tabular}
}
\end{table*}

Refinement (discrimination) decreases substantially in instruction-tuned models: by 0.0216 in Qwen (95\% CI: [0.0179, 0.0253]) and 0.0487 in Mistral (95\% CI: [0.0453, 0.0522]). Uncertainty remains approximately constant, as expected since it depends only on base rates. Reliability increases (worse calibration) in both families, particularly Mistral.

The refinement degradation is consistent with geometric distortion affecting the shape of the probability landscape, rather than scalar distortion that would primarily affect reliability.

\subsubsection{5.5 Summary of Hypothesis Tests}

\textbf{Table 5: Hypothesis Gate Summary}

\begin{table*}[htbp]
\centering
\resizebox{\dimexpr 2\columnwidth+\columnsep\relax}{!}{%
\begin{tabular}{lllll}
\toprule
Hypothesis & Type & Metric & Result & Gate \\
\midrule
H-E1 & EXISTENCE & AUROC degradation & +0.0303 mean & PASS \\
H-M1 & MECHANISM & Margin inflation & 9.92x mean & PASS \\
H-M2 & MECHANISM & $\beta$ attenuation & 0.69 mean & PASS \\
H-M3 & MECHANISM & Refinement delta & -0.035 mean & PASS \\
\bottomrule
\end{tabular}
}
\end{table*}

All four sub-hypotheses show results in the predicted direction with statistical significance.
